%!TEX TS-program = xelatex
\documentclass[12pt,a4paper]{extarticle} 

\input{misc/_preamble.tex}
\addbibresource{misc/_references.bib}

\def\course{Теория игр}
\def\groupnumber{$\gamma$}
\date{7 августа 2026 года}
\def\lecturer{Сергей Кыльчик}
\def\assist{Климент Шамаев}
\def\listokname{Модель Хотеллинга}
\boolfalse{onlyorder}

\begin{document}
\input{misc/_lesh2026.tex}

%-------------------------------------------------------------------------
\z[Классика]{
  Город $N$ представляет собой отрезок длины $1$, на котором равномерно распределены $120$ жителей. На краях города расположены фирмы. Функция издержек первой фирмы имеет вид $TC_1=20q_1$, функция издержек второй фирмы --- $TC_2=cq_2$. Каждый из жителей города покупает ровно единицу товара. Транспортные издержки равны расстоянию. Фирмы конкурируют, одновременно и независимо выбирая цены.
}{
  \n а) Найдите все рыночные равновесия для $c=20$.
  \n \textbf{Решение:}
  Спрос на продукцию фирмы 1 (левой):
  \[
  D_1(p_1,p_2)=
  \begin{cases}
    0, & p_1-p_2\ge 1,\\[4pt]
    \dfrac{1+p_2-p_1}{2}, & -1 < p_1-p_2 < 1,\\[4pt]
    1, & p_1-p_2\le -1.
  \end{cases}
  \]
  $D_2=1-D_1$. При $c=20$ издержки равны. Кривые реакции (наилучшие ответы):
  \[
  BR_1(p_2)=
  \begin{cases}
    [p_2+1,\infty), & p_2<19,\\[4pt]
    \left\{\dfrac{21+p_2}{2}\right\}, & 19\le p_2\le 23,\\[4pt]
    \{p_2-1\}, & p_2>23,
  \end{cases}
  \qquad
  BR_2(p_1)=
  \begin{cases}
    [p_1+1,\infty), & p_1<19,\\[4pt]
    \left\{\dfrac{21+p_1}{2}\right\}, & 19\le p_1\le 23,\\[4pt]
    \{p_1-1\}, & p_1>23.
  \end{cases}
  \]
  Пересечение даёт единственную точку $(21,21)$. Следовательно, равновесие:
  \[
  \boxed{p_1=p_2=21,\quad q_1=q_2=60.}
  \]

  \n б) Найдите все рыночные равновесия в зависимости от $c$.
  \n \textbf{Решение:}
  Общие кривые реакции:
  \[
  BR_1(p_2)=
  \begin{cases}
    [p_2+1,\infty), & p_2<19,\\[4pt]
    \left\{\dfrac{21+p_2}{2}\right\}, & 19\le p_2\le 23,\\[4pt]
    \{p_2-1\}, & p_2>23,
  \end{cases}
  \]
  \[
  BR_2(p_1)=
  \begin{cases}
    [p_1+1,\infty), & p_1<c-1,\\[4pt]
    \left\{\dfrac{1+p_1+c}{2}\right\}, & c-1\le p_1\le c+3,\\[4pt]
    \{p_1-1\}, & p_1>c+3.
  \end{cases}
  \]
  Найдём все пары $(p_1,p_2)$, удовлетворяющие $p_1\in BR_1(p_2)$ и $p_2\in BR_2(p_1)$. Перебор всех комбинаций кусков даёт следующие множества:

  \begin{itemize}
    \item \textbf{Случай $c<17$:}
    \[
    \boxed{
      p_2 = p_1 - 1,\qquad p_1 \in [c+3,\; 20].
    }
    \]
    \item \textbf{Случай $17\le c\le 23$:}
    \[
    \boxed{
      p_1=\frac{43+c}{3},\qquad p_2=\frac{23+2c}{3}.
    }
    \]
    \item \textbf{Случай $c>23$:}
    \[
    \boxed{
      p_2 = p_1 - 1,\qquad p_2 \in [23,\; c] \quad(\text{т.е. } p_1\in[22,\; c-1]).
    }
    \]
  \end{itemize}
  Других решений нет, что проверяется подстановкой во все оставшиеся комбинации (они дают либо противоречия, либо уже включены в указанные множества).
}


%-------------------------------------------------------------------------
\z[Это проще, чем вы думаете]{
  Потребители некоторого товара живут на отрезке $[0,1]$. Потребители распределены равномерно. Фирмы-производители расположены на краях города. Фирмы одновременно и независимо выбирают цены. Издержки каждой фирмы: $TC=3q$. Каждый потребитель покупает ровно единицу товара, минимизируя суммарные издержки. Транспортные издержки равны расстоянию.
}{
  \n а) Найдите спрос на продукцию каждой из фирм для любых $P_1,P_2\ge 0$.
  \n \textbf{Решение:}
  Точка безразличия: $x = (1+P_2-P_1)/2$. Тогда
  \[
  D_1(P_1,P_2)=
  \begin{cases}
    0, & P_1-P_2\ge 1,\\[4pt]
    \dfrac{1+P_2-P_1}{2}, & -1 < P_1-P_2 < 1,\\[4pt]
    1, & P_1-P_2\le -1,
  \end{cases}
  \qquad
  D_2=1-D_1.
  \]

  \n б) Найдите рыночное равновесие.
  \n \textbf{Решение:}
  Выведем кривые реакции. Для фирмы 1 при заданном $P_2$:
  \[
  BR_1(P_2)=
  \begin{cases}
    [P_2+1,\infty), & P_2<2,\\[4pt]
    \left\{2+\dfrac{P_2}{2}\right\}, & 2\le P_2\le 6,\\[4pt]
    \{P_2-1\}, & P_2>6.
  \end{cases}
  \]
  Аналогично для фирмы 2:
  \[
  BR_2(P_1)=
  \begin{cases}
    [P_1+1,\infty), & P_1<2,\\[4pt]
    \left\{2+\dfrac{P_1}{2}\right\}, & 2\le P_1\le 6,\\[4pt]
    \{P_1-1\}, & P_1>6.
  \end{cases}
  \]
  Пересекаем множества. Единственное решение в средних частях:
  \[
  P_1 = 2+\frac{P_2}{2},\qquad P_2 = 2+\frac{P_1}{2}
  \;\Longrightarrow\; P_1=P_2=4.
  \]
  Другие комбинации (левая/средняя, правая/средняя и т.д.) не дают пересечений (проверяется подстановкой). Следовательно, единственное равновесие:
  \[
  \boxed{P_1=P_2=4}.
  \]

  \n в) Государство вводит потолок цен на первую фирму $\overline{P_1}$ и на вторую $\overline{P_2}$. Найдите новое рыночное равновесие.
  \n \textbf{Решение:}
  Теперь цены ограничены: $P_1\le \overline{P_1}$, $P_2\le \overline{P_2}$. Кривые реакции с учётом ограничений:
  \[
  BR_1(P_2;\overline{P_1}) =
  \begin{cases}
    \min\{\overline{P_1},\; P_2+1\}, & P_2<2,\\[4pt]
    \min\{\overline{P_1},\; 2+\frac{P_2}{2}\}, & 2\le P_2\le 6,\\[4pt]
    \min\{\overline{P_1},\; P_2-1\}, & P_2>6,
  \end{cases}
  \]
  причём если $\min$ даёт цену ниже 3, то прибыль отрицательна, и фирма предпочтёт не продавать, т.е. выберет $P_1 = \min\{\overline{P_1}, P_2+1\}$ (при этом спрос равен 0, прибыль 0). Аналогично для фирмы 2.

  Равновесие находится как решение системы:
  \[
  P_1 = BR_1(P_2;\overline{P_1}),\qquad P_2 = BR_2(P_1;\overline{P_2}).
  \]
  Анализ этой системы даёт следующие случаи:
  \begin{itemize}
    \item Если $\overline{P_1}\ge 4$ и $\overline{P_2}\ge 4$, то равновесие $(4,4)$.
    \item Если $\overline{P_1}<4$ и $\overline{P_2}\ge 4$, то равновесие $P_1=\overline{P_1},\; P_2=2+\frac{\overline{P_1}}{2}$ (при условии, что $P_2\le \overline{P_2}$, иначе переходим к следующему случаю).
    \item Если $\overline{P_2}<4$ и $\overline{P_1}\ge 4$, то симметрично: $P_2=\overline{P_2},\; P_1=2+\frac{\overline{P_2}}{2}$.
    \item Если оба потолка меньше 4, то равновесие может быть только если один из потолков «связывает», а другой нет. В общем виде ответ задаётся системой выше.
  \end{itemize}
  Таким образом, $\boxed{\text{равновесие определяется системой кривых реакций с ограничениями}}$.

  \n г) Найдите все общественно-оптимальные значения $\overline{P_1},\overline{P_2}$.
  \n \textbf{Решение:}
  Полезность потребителей от потребления и издержки фирм на производство фиксированны. Поэтому надо оптимизировать только транспортные издержки, они имеют вид $0.5x^2+0.5(1-x)^2$. Получаем $x=0.5$, достигается при $\overline{P_1},\overline{P_2}\ge 4$ или $\overline{P_1}=\overline{P_2}$
}

%-------------------------------------------------------------------------
\z[Это сложнее, чем вы думаете]{
  Город состоит из одной улицы длиной 1 км, на которой равномерно расположено население. В концах города расположены две фирмы: А и Б. Издержки фирм равны $0$. Фирмы одновременно назначают цены. Транспортные издержки: на половине улицы, ближайшей к фирме Б, издержки в 2 раза выше.
}{
  \item Просто все как обычно.
  \item Та половина улицы, которая ближе к фирме Б, оказалась затоплена: издержки перемещения по ней оказываются в 2 раза выше.
}
$$\includegraphics[scale=0.4]{Подборки/Микроэкономика/solv1.jpg}$$
$$\includegraphics[scale=0.4]{Подборки/Микроэкономика/solv7.jpg}$$
$$\includegraphics[scale=0.4]{Подборки/Микроэкономика/solv5.jpg}$$
$$\includegraphics[scale=0.4]{Подборки/Микроэкономика/solv2.jpg}$$
$$\includegraphics[scale=0.4]{Подборки/Микроэкономика/solv6.jpg}$$
$$\includegraphics[scale=0.4]{Подборки/Микроэкономика/solv4.jpg}$$
$$\includegraphics[scale=0.4]{Подборки/Микроэкономика/solv3.jpg}$$
$$\includegraphics[scale=0.4]{Подборки/Микроэкономика/solv8.jpg}$$
Пытаемся пересечь, получаем, что равновесия нет.
}

\end{document}